\section{数据预处理与泄漏控制}

\subsection{跨年度数据审计}

正式分析前，先对2015--2024年共10年的原始数据进行系统性审计：

\begin{table}[H]
\centering
\caption{跨年度数据审计内容}
\label{tab:audit}
\begin{tabular}{p{3.5cm}p{8.5cm}}
\hline
\textbf{审计项目} & \textbf{检查内容与输出} \\
\hline
字段名称一致性 & 各年份CSV的列名是否变化（如大小写、空格、括号差异），生成字段映射表 \\
\hline
字段单位一致性 & 同一字段在不同年份的单位是否一致（如kBtu vs GJ），生成字段映射表中的单位映射 \\
\hline
建筑类型命名一致性 & BuildingType和LargestPropertyUseType的取值在不同年份是否统一，是否存在拼写差异或分类合并 \\
\hline
建筑ID稳定性 & OSEBuildingID是否跨年稳定，是否存在同一ID在不同年份指向不同建筑的情况 \\
\hline
重复记录检查 & 同一建筑在同一年份是否有多条记录，完全重复行比例 \\
\hline
数值字段解析 & 是否存在因格式问题（如千分位逗号、特殊字符）导致数值字段被解析为字符串的情况 \\
\hline
各字段覆盖率 & 每年各字段的非缺失比例，生成字段覆盖率统计表 \\
\hline
PropertyGFATotal核查 & 重新检查该字段在各年份的实际覆盖率，确认是否因字段名称变化导致误判为全部缺失 \\
\hline
\end{tabular}
\end{table}

\subsection{数据预处理流程}

数据预处理按以下顺序执行，所有拟合操作（填充值、编码映射、标准化参数）仅在训练集上计算：

\begin{table}[H]
\centering
\caption{预处理步骤说明}
\label{tab:preprocessing}
\begin{tabular}{p{1cm}p{3.5cm}p{7.5cm}}
\hline
\textbf{步骤} & \textbf{操作} & \textbf{说明} \\
\hline
P1 & 字段统一 & 去除列名中括号、空格、斜杠等特殊字符；统一大小写；修正各年份字段名差异 \\
\hline
P2 & 数值类型转换 & 将应为数值型的列强制转换为numeric，无法解析的值标记为NaN \\
\hline
P3 & 唯一建筑与重复检查 & 同年同OSEBuildingID重复、完全重复行检查；记录重复条数与处理方式 \\
\hline
P4 & 目标泄漏审查 & 识别并剔除与SiteEUI存在确定性数量关系或共享年度能耗测量信息的派生字段（SourceEUI、SiteEnergyUse、SourceEnergyUse、天气标准化字段等） \\
\hline
P5 & 目标变量过滤 & 仅删除SiteEUI缺失、SiteEUI $\le$ 0和明确录入错误的记录；\textbf{不进行分位数截尾}，保留合理的极端高能耗建筑作为后续分析对象 \\
\hline
P6 & YearBuilt处理 & YearBuilt $>$ DataYear的异常值标记为缺失，增加异常指示变量YearBuilt\_Anomaly；按建筑类型中位年份填补；确保BuildingAge非负 \\
\hline
P7 & 缺失值填充 & 数值字段优先以同建筑类型中位数填充，仍缺失时使用训练集整体中位数；类别字段填充为``Unknown''；同时生成各字段的缺失指示变量（如ENERGYSTARScore\_missing） \\
\hline
P8 & 特征工程 & 计算BuildingAge、LogGFA、面积比率、混合用途标识、缺失指示变量等（详见第5节特征列表） \\
\hline
P9 & 类别编码 & ElasticNet使用OneHotEncoder；树模型使用OneHotEncoder（LightGBM可选用原生类别特征）；\textbf{禁止对建筑类型使用LabelEncoder}以避免虚假的大小关系 \\
\hline
P10 & 目标变换与标准化 & 建模时对目标取$y = \log(1+\text{SiteEUI})$，预测后还原$\hat{y}_{\text{orig}} = \exp(\hat{y}_{\text{log}})-1$；RobustScaler在训练集上拟合，对验证集和测试集仅transform \\
\hline
\end{tabular}
\end{table}

\subsection{目标泄漏审查}

目标泄漏是建筑能耗建模中最关键的风险控制措施。在预测SiteEUI时，任何从同期能耗信息中派生的字段若进入特征集，都将导致模型在训练时``看到''目标的部分信息，从而高估真实预测能力。本课题在建模前对所有字段逐一审查，确保以下三类问题不会发生：

\textbf{第一类：直接泄漏——由目标变量数学变换得到的字段。}SourceEUI由SiteEUI乘以源-站转换系数计算得出（$\text{SourceEUI} = \text{SiteEUI} \times \text{Source-to-Site Factor}$）；SiteEnergyUse为建筑面积 $\times$ SiteEUI。这些字段与目标存在确定性的数量关系，一经发现立即剔除。

\textbf{第二类：间接泄漏——与目标共享年度能耗测量信息的字段。}SourceEnergyUse、SiteEnergyUseWN、SourceEUIWN等字段虽然计算方式不同，但均来自同一份年度能耗账单数据。在预测场景下（使用静态属性估计合理能耗水平），这些字段的值无法提前获知，因此必须剔除。

\textbf{第三类：灰色变量——与目标相关但非直接派生的字段。}TotalGHGEmissions和GHGEmissionsIntensity由能耗乘以排放因子计算，预测模型中不可使用，但在后续的优先级排序阶段可作为减碳维度的参考指标。ENERGYSTARScore实际由EPA基于当年SiteEUI百分位计算，具有较强的目标相关性，因此主模型（模型A）不使用，仅编码为二值指示变量（ES\_Missing）；模型B使用ENERGYSTARScore仅作为扩展对比实验。

\begin{table}[H]
\centering
\caption{目标泄漏审查与字段处理决策}
\label{tab:leakage}
\begin{tabular}{p{3.5cm}p{2.5cm}p{6cm}}
\hline
\textbf{字段} & \textbf{处理} & \textbf{理由} \\
\hline
SourceEUI、SourceEUIWN & 剔除 & 由SiteEUI乘以源-站转换系数计算，与目标存在确定性数量关系 \\
\hline
SiteEnergyUse、SiteEnergyUseWN、SourceEnergyUse & 剔除 & 能耗总量 $\div$ 面积可直接还原为SiteEUI，属直接泄漏 \\
\hline
TotalGHGEmissions、GHGEmissionsIntensity & 预测模型剔除，排序可用 & 由能耗和排放因子计算，预测时不可预先获得；排序阶段作为减碳参考 \\
\hline
ENERGYSTARScore & 主模型不直接入模，仅保留ES\_Missing指示变量；模型B使用作扩展对比 & 由EPA基于当年SiteEUI百分位计算，与目标高度相关；防止模型因评级信息而虚高 \\
\hline
PropertyGFATotal & 需重新核查各年份实际覆盖率 & 不应直接假定全部缺失，可能因字段名称变化导致匹配失败 \\
\hline
\end{tabular}
\end{table}

\subsection{SiteEUI异常值处理策略}

高能耗建筑正是本研究的识别对象，因此\textbf{不能对SiteEUI简单分位数截尾}。处理策略为：

\begin{enumerate}
    \item 删除SiteEUI缺失、非正值和明确录入错误（如负数、数量级异常的错误值）的记录；
    \item \textbf{保留合理的极端高能耗记录}，它们将在后续的异常识别和优先级排序中被标记和分析；
    \item 建模时对目标取对数$y = \log(1+\text{SiteEUI})$以缓解右偏分布对模型训练的影响；
    \item 预测后通过$\hat{y}_{\text{orig}} = \exp(\hat{y}_{\text{log}}) - 1$恢复为原始尺度；
    \item 原始SiteEUI继续用于异常识别和优先级排序。
\end{enumerate}

\subsection{建筑类型—年份基准匹配}

利用基准性能范围数据（Benchmarking Performance Ranges by Building Type）：

\begin{enumerate}
    \item 为每栋建筑匹配同类型、同年度的SiteEUI基准值（如P50中位数和P75上四分位数）；
    \item 计算超耗比例：正值表示高于同类平均水平，负值表示低于平均水平；
    \item 在优先级评分阶段，使用同类型超耗比例替代绝对SiteEUI，降低建筑类型固有差异的干扰，避免办公楼、医院等类型因天然能耗高而垄断Top-50。
\end{enumerate}

\subsection{目标变换}

SiteEUI通常呈右偏分布。本课题在训练时对目标进行$\log(1+\text{SiteEUI})$变换以改善残差分布的对称性，评估时在原始尺度（kBtu/sf）上报告误差。

\subsection{划分策略}

采用按年份的扩展窗口滚动验证（详见第5节），不再使用TimeSeriesSplit的机械划分。同一建筑可能出现在多个集合中（因为同栋建筑在不同年份均有记录），这符合``利用历史数据估计基准''的场景设定。辅助实验采用按OSEBuildingID的GroupShuffleSplit，检验模型对完全未见建筑的泛化能力。

\subsection{技术可行性}

本研究的技术基础坚实：

\begin{enumerate}
    \item \textbf{数据可用性：}西雅图市政府公开发布，数据完整规范，附带元数据说明，覆盖2015--2024共10年；
    \item \textbf{工具链成熟度：}Python Scikit-learn、XGBoost、LightGBM、SHAP均为工业级稳定库；
    \item \textbf{计算资源：}数据规模约15 MB，在合理控制参数搜索范围的情况下，可在课程项目周期内完成全部训练与解释分析；
    \item \textbf{方法论支撑：}回归预测、SHAP可解释性、K-Means聚类、多维加权评分均为经过充分验证的方法，本研究在此基础上进行整合与改进；
    \item \textbf{可复现性：}全流程Python脚本化，环境依赖写入requirements.txt；
    \item \textbf{经验基础：}已完成鸢尾花数据集分析项目，具备完整的ML pipeline实施经验。
\end{enumerate}

\subsection{主要难点与应对}

\begin{table}[H]
\centering
\caption{技术难点与应对策略}
\label{tab:challenges}
\begin{tabular}{p{3cm}p{4cm}p{5cm}}
\hline
\textbf{难点} & \textbf{原因} & \textbf{应对策略} \\
\hline
缺失值处理 & 部分字段缺失率达20--40\% & 数值以同建筑类型中位数优先填充，类别填充``Unknown''，同时生成缺失指示变量保留缺失信息 \\
\hline
不同建筑类型样本量差异大 & 办公楼远多于医院等类型 & 评估时按建筑类型分组报告性能，排序使用同类型超耗比例避免类型垄断 \\
\hline
目标泄漏风险 & 派生变量与目标共享能耗信息 & 严格剔除同期派生字段，仅保留预测时可获得的静态属性 \\
\hline
信息泄漏于测试集 & 缺失值填充等操作在全量数据执行 & 所有拟合（中位数、编码器、标准化器）仅在训练集上计算 \\
\hline
高能耗建筑极端值 & 无法区分录入错误与真实高耗能 & 采用log1p变换缓解偏度；结合预测残差和同类基准双重验证 \\
\hline
K值选择主观 & 不同K值给出不同分群结果 & 综合轮廓系数+Calinski-Harabasz指数+画像可解释性综合判断 \\
\hline
权重系数主观 & 优先级评分中各维度重要性缺乏客观标准 & 设计多个权重方案（节能优先/减碳优先/大型建筑优先），通过Bootstrap、Jaccard相似度、Kendall相关系数检验排序稳定性 \\
\hline
PropertyGFATotal高缺失 & 可能因字段名称变化导致匹配失败 & 跨年度审计时重新核查，若确认为真实缺失则使用PropertyGFABuildings替代 \\
\hline
\end{tabular}
\end{table}
